% Shared module: identification, uniqueness, the spectral form, limits of certification.
\subsection{When does a solution determine its equation?}\label{sec:ident}
Let $\Omega$ be open and connected (on a disconnected set the pieces may obey different equations),
let $u\in\mathcal D'(\Omega)$ satisfy $P(\partial)u=0$, and let $P=\prod_iP_i^{e_i}$ (up to a nonzero constant factor) be the
factorization of $P$ into pairwise non-associate irreducible polynomials $P_i$ with multiplicities
$e_i$. The question is whether $\Ann(u)=(P)$: then every constant-coefficient equation satisfied by
$u$ is a consequence of $P$, and $P$ is determined by $u$.

\begin{theorem}[Identification criterion]\label{thm:ident}
Put $v_i=(P/P_i)(\partial)\,u$ and $v_i'=(P/P_i^{e_i})(\partial)\,u$.
\begin{enumerate}[label=(\alph*)]
  \item $\Ann(u)=(P)$ if and only if $v_i\notin\EP(\Omega)$ for every $i$.
  \item $\Var(\Ann u)=\Var(P)$ if and only if $v_i'\notin\EP(\Omega)$ for every $i$.
  \item If $P$ is square-free, (a) and (b) coincide. If $P$ is monic in $\tau$ of degree $m$, then
  $\Ann(u)=(P)$ if and only if the only $(\varphi_0,\dots,\varphi_{m-1})\in\C[\xi]^m$ with
  $\sum_{j<m}\varphi_j(\partial_x)\,\partial_t^{\,j}u=0$ is zero.
\end{enumerate}
\end{theorem}

The idea is to strip one irreducible factor and look at what is left. The function $v_i$ solves
$P_i(\partial)v_i=0$. If $v_i$ is a generic solution of that factor, no smaller operator kills $u$.
If $v_i$ is an exponential-polynomial, some $h$ coprime to $P_i$ kills it, and then
$h\cdot P/P_i$ is an annihilator of $u$ that is not a multiple of $P$. The converse uses
\cref{lem:gcd}; the full proof is in \cref{app:proofs}.

\begin{proposition}[Uniqueness and the intrinsic operator]\label{prop:unique}
\begin{enumerate}[label=(\alph*)]
  \item If $\Ann(u)=(P)$, then $P$ is determined by $u$ up to a nonzero constant and is the unique
  annihilator of minimal degree.
  \item $\Ann(u)=(P)$ if and only if $\sigp(u)=\emptyset$ and $\Gu\sim P$ (equality up to a constant).
  \item For any annihilator $P$, the multiplicity of $P_i$ in $\Gu$ equals $e_i-\max\{k\le e_i:
  (P/P_i^{k})(\partial)u\in\EP\}$. Hence $\Gu$ is computed from any annihilator by at most
  $\sum_ie_i$ membership tests ``$\in\EP$''.
\end{enumerate}
\end{proposition}

\begin{example}[The heat equation, continued]\label{ex:heat-ident}
For $P=\tau-\nu\xi^2$ (irreducible) the criterion reads: $\Ann(u)=(P)$ if and only if $u\notin\EP$.
One mode $e^{-\nu\kappa^2t+i\kappa x}$ has $\Ann(u)=(\tau+\nu\kappa^2,\ \xi-i\kappa)$, a point of
$\Var(P)$; it cannot distinguish the heat equation from the pure decay $\partial_t+\nu\kappa^2$, or from
any other equation whose variety passes through the point $(-\nu\kappa^2,i\kappa)$. The heat kernel is not in $\EP$, so every
constant-coefficient equation it satisfies is a multiple of $\partial_t-\nu\partial_x^2$.
\end{example}

\begin{remark}[Curves and points]\label{rem:curves-points}
The two cases of \cref{ex:heat-ident} sit at the two ends of the decomposition~\eqref{eq:GJ}. A finite sum of modes
$u\in\EP$ has an annihilator with finitely many zeros, so
\[
  u\in\EP:\qquad \Gu=1,\qquad \sigp(u)=\{(\tau_j,\xi_j)\}\ \text{finite (its exponents)},
\]
whereas the zero set of a nonconstant polynomial in two variables is a whole algebraic curve. For the heat kernel
the situation is reversed:
\[
  u=t^{-1/2}e^{-x^2/(4\nu t)}:\qquad \Gu=\tau-\nu\xi^2,\quad \Var(\Gu)\ \text{a curve},\qquad \sigp(u)=\emptyset .
\]
By \cref{prop:unique}(b) it is exactly this pair, an empty point spectrum together with $\Gu\sim P$, that gives
$\Ann(u)=(P)$: the solution fills the whole curve of its equation and contains no isolated mode besides.
\end{remark}

\roleinalg{Data never contain the heat kernel, they contain finitely many sampled modes; a finite
set of modes is always an exponential-polynomial. The algorithm therefore never claims that an
equation is \emph{identified}. It asks whether the modes it sees are \emph{compatible} with a factored
symbol and are rich enough in each factor (enough distinct wavenumbers per branch) for the factor to
be pinned down within its class. This is the meaning of the label \FC{} in the companion algorithm paper.}

\subsection{What finite data can and cannot certify}\label{sec:certify}
\Cref{thm:ident} speaks about $u$, not about finitely many numbers. Two facts frame what an algorithm
can promise.

Recall that $\ring u=\spn\{\partial_t^{\,j}\partial_x^{\,k}u:\ j,k\ge0\}$ is the vector space generated by the
derivatives of $u$ (\cref{sec:ep}). The \emph{jet} of order $M$ of $u$ at a point $z_0=(t_0,x_0)$ is the list of its
derivatives of order at most $M$ at that point; for example
\[
  j^2_{z_0}(u)=\big(u,\ u_t,\ u_x,\ u_{tt},\ u_{tx},\ u_{xx}\big)\big|_{(t,x)=z_0}.
\]

\begin{proposition}[Finite data]\label{prop:certify}
Let $\Omega$ be connected and $z_0\in\Omega$.
\begin{enumerate}[label=(\roman*)]
  \item If $\dim_\C\ring u\le n$, the jet of order $2n-1$ of $u$ at $z_0$ determines $\Ann(u)$ among all functions
  $u'$ with $\dim_\C\ring u'\le n$.
  \item For every $u$ and every $M$ there is a polynomial, hence an exponential-polynomial, with the same jet of
  order $M$ at $z_0$; and for finitely many sample points there is an exponential-polynomial with the same values.
  Consequently, without a bound on $\dim\ring u$, no finite set of jets or samples proves $\Ann(u)=(P)$.
\end{enumerate}
\end{proposition}

\begin{proof}
(ii) The Taylor polynomial of order $M$ at $z_0$ has the same jet, and polynomial interpolation in two variables
matches any finite set of values. Both are exponential-polynomials, and an exponential-polynomial never has
$\Ann=(P)$ when $P$ has an irreducible factor (\cref{thm:ident}: all $v_i$ are then exponential-polynomials).
(i) The space $\ring u$ is invariant under differentiation and, by \cref{lem:ep}, consists of real-analytic
exponential-polynomials. Let $K_M$ be the subspace of elements of $\ring u$ whose jet of order $M$ vanishes at $z_0$.
Then $K_{M+1}=\{f\in K_0:\ \partial_tf,\partial_xf\in K_M\}$, so once $K_{M+1}=K_M$ the chain is constant, and a
constant $K_M$ is a derivative-invariant space of functions vanishing at $z_0$ with all their derivatives, hence
$\{0\}$. Since $\dim K_0\le n-1$, the chain reaches $\{0\}$ by $M=n-1$: the jet of order $n-1$ at $z_0$ is injective
on $\ring u$, and likewise on $\ring u'$. For $Q$ of degree at most $n$, $Q\in\Ann(u)$ iff the jet of order
$n-1$ of $Q(\partial)u$ vanishes at $z_0$, a condition on the jet of order $2n-1$ of $u$; the same holds for $u'$.
Finally $\Ann(u)$ has codimension at most $n$, so a Gr\"obner basis for a degree-compatible order consists of
polynomials of degree at most $n$; the same holds for $\Ann(u')$, and two ideals generated in degree at most $n$
that contain the same polynomials of degree at most $n$ are equal.
\end{proof}

The consequence for learning is a design rule: without an a priori bound on $\dim\ring u$, which data never
supply, an algorithm can refute a candidate equation, or report that it is compatible with the data at
the resolution available, but it cannot certify it. Its output must carry a label that says which of these
happened.
